% Appendix A-C. Numbers from RESULTS.md ("Implementation checks", "Serving-cost roofline model",
% "Simulations of next methods", "Tokenizer level-skip probe").

\section{Implementation checks}
\label{sec:checks}

\begin{itemize}
  \item With $M{=}1$ the chunked model reproduces the released model's teacher-forced loss for $G{=}0$ and $G{=}11$
        (1.290056 vs 1.290056).
  \item The tokenizer round trip on real 128-token chunks gives 100\% token accuracy.
  \item The per-level training path (one short sequence per level) equals the single dense-mask sequence for
        $M{=}4$ and $G\in\{0,1,3,6,11\}$ to within $10^{-6}$; the dense path does not fit in 8\,GB.
  \item The KV-cached plan gives the same loss as recomputing the plan at every local level (difference
        $\le 8{\times}10^{-7}$ for $G\in\{0,1,3,6,11\}$) and the same generated codes under the same seed; it halves
        the time of a training step (1.21 vs 2.43\,s on a T4).
  \item Halo with $M{=}1$ changes nothing. Training inputs equal generation inputs bit for bit at every level.
\end{itemize}

\section{Serving-cost model}
\label{sec:cost}

The cost model (\texttt{paper/scripts/serving\_cost.py} and \texttt{method\_scenarios.py}) is analytic, not measured.
It assumes 7B parameters for both an AR model (grouped-query attention with 8 KV groups) and NCM, H100 specifications
(989 TFLOP/s, 3.35\,TB/s, 80\,GB), 40\% FLOP and 80\% bandwidth utilisation, \$2.50 per GPU-hour, and 2,048-token
outputs ($M{=}16$). It counts dense and attention FLOPs, weight loads and AR KV-cache reads, and ignores the tokenizer's
decoder, kernel-launch overhead and low utilisation on small passes. It assumes equal quality at equal size, which
this paper does not show.

\begin{table}[h]
\centering
\small
\caption{Cost model, 7B parameters on an H100, 2,048-token outputs.}
\begin{tabular}{@{}lcccc@{}}
\toprule
method & positions per token & batch-1 latency & vs AR & AR cheaper from batch \\
\midrule
AR (GQA-8) & 1 & 10.75\,s & 1$\times$ & -- \\
B1, chunks in sequence & 5.00 & 0.92\,s & 12$\times$ & 35 \\
B2, fully global & 5.00 & 0.38\,s & 28$\times$ & 33 \\
Block-NCM $G{=}3$, $h{=}4$ & 5.50 & 0.41\,s & 26$\times$ & 31 \\
Block-NCM $G{=}5$, $h{=}4$ & 5.38 & 0.40\,s & 27$\times$ & 31 \\
\bottomrule
\end{tabular}
\end{table}

At batch 1 every NCM variant is limited by loading weights, and B1 is about 2.3$\times$ slower than Block-NCM
with halo because it makes $M$ times as many passes, each over fewer positions. At large batch NCM is compute-bound: it processes about 5 positions per
output token, so an AR model becomes cheaper per token from about batch 31. Without grouped-query attention the AR
model's KV cache (1.07\,GB per sequence) caps it below batch 64 on 80\,GB, and Block-NCM is cheaper at every feasible
batch. Block-NCM and B2 differ little at this size: attention is a small share of FLOPs at $d{=}4096$, and the
saving from local attention matters only for much longer outputs.

\section{Tokenizer level-skip probe}
\label{sec:probe}

Can the frozen residual tokenizer reconstruct text when some levels are never produced? We encode 1,024 real 128-token
chunks, drop levels, decode, and compare tokens. \emph{Re-quantised}: the residual quantiser is re-run without those
levels, so later levels absorb the residual (the best case, for a generator trained on that tokenisation).
\emph{Hole}: the normal codes with the skipped levels left out at decode time.

\begin{table}[h]
\centering
\small
\caption{Token accuracy and exactly reconstructed chunks when levels are skipped (levels numbered 1--11).}
\begin{tabular}{@{}lccccc@{}}
\toprule
skipped levels & positions per chunk & re-quantised: tokens & exact chunks & hole: tokens & exact chunks \\
\midrule
none & 641 & 99.97\% & 96.7\% & 99.97\% & 96.7\% \\
11 (the last) & 513 & 98.26\% & 17.9\% & 98.26\% & 17.9\% \\
7 & 577 & 99.90\% & 88.7\% & 99.61\% & 62.7\% \\
7, 8, 10 (keeps 6$\to$9$\to$11) & 393 & 99.52\% & 57.0\% & 98.41\% & 14.4\% \\
7--10 (keeps 6$\to$11) & 301 & 98.71\% & 23.5\% & 97.10\% & 3.6\% \\
4--10 & 177 & 91.89\% & 0\% & 84.81\% & 0\% \\
\bottomrule
\end{tabular}
\end{table}

Intermediate levels can be skipped with little loss when the codes are re-quantised: keeping levels 6, 9 and 11
preserves 99.5\% of tokens with 39\% fewer positions. The last level cannot be skipped. These are upper bounds with
perfect codes; a generator would still have to predict the larger jumps.
